% from ega4-8-fr.tex
\oldpage[IV-3]{42}
point quelconque de $Y$, donc on voit finalement que l'on peut supposer que
$Y=\operatorname{Spec}(A)$, où $A$ est un \emph{anneau local noethérien}.
Soient $\mathfrak m$ l'idéal maximal de $A$, $\widehat A$ le complété de $A$
pour la topologie $\mathfrak m$-préadique~; on sait que $\widehat A$ est un
anneau local noethérien et que le morphisme
$\operatorname{Spec}(\widehat A)\to\operatorname{Spec}(A)$ est fidèlement
plat et quasi-compact (0\textsubscript{I}, 7.3.5)~; quant à (2.7.1, (i),
(vii), (xiv), (xv) et (xvi)), on voit en outre qu'on peut se borner au cas où
$A$ est \emph{complet}. Il résulte alors de (II, 6.2.6) que
$X=X'\amalg X''$, où $X'$ est un $Y$-schéma fini et $X''$ un $Y$-schéma
quasi-fini tel que $X''\cap f^{-1}(y)=\varnothing$.
\end{env}

Plaçons-nous d'abord dans les hypothèses de (8.11.1)~; comme $f$ est propre,
$X''$, qui est fermé dans $X$, est propre sur $Y$ (II, 5.4.10), donc
$f(X'')$ est fermé dans $Y$~; mais $y$ n'est pas contenu dans $f(X'')$, et
par ailleurs est adhérent à tout point de $Y$, donc $f(X'')=\varnothing$, et
par suite $X''$ est vide et $f$ est \emph{fini}.

Plaçons-nous maintenant dans les hypothèses de (8.11.2) et, en nous ramenant
encore (comme on peut le faire d'après ce qui précède) au cas où
$Y=\operatorname{Spec}(A)$ est affine et noethérien de dimension finie,
raisonnons en outre par récurrence sur la dimension de $Y$. En se ramenant
comme ci-dessus au cas où $A$ est en outre local et complet, on a
$\dim(\mathcal O_y)=\dim(A)=\dim(Y)$ et pour tout $z\neq y$,
$\dim(\mathcal O_z)<\dim(\mathcal O_y)$, donc
$\dim(Y-\{y\})<\dim(Y)$. Or, par hypothèse on a
$f(X'')\subset Y-\{y\}$ et la restriction de $f$ à $X''$ est évidemment un
morphisme quasi-fini et séparé~; appliquant à $Y-\{y\}$ et à $X''$
l'hypothèse de récurrence, on voit que $X''$ est quasi-affine sur
$Y-\{y\}$~; mais $Y-\{y\}$ étant quasi-affine sur $Y$ puisque $Y$ est
noethérien, $X''$ est aussi quasi-affine sur $Y$ (II, 5.1.10, (ii))~; comme
par ailleurs $X'$ est fini (et \emph{a fortiori} affine) sur $Y$, $X$ est
quasi-affine sur $Y$ (II, 4.6.17 et 5.1.2, c')).

\begin{proposition}[8.11.5]
\label{IV.8.11.5-fr}
Soit $f:X\to Y$ un morphisme de présentation finie. Les conditions suivantes
sont équivalentes~:
\begin{enumerate}
\item[\emph{a)}] $f$ est une immersion fermée.
\item[\emph{b)}] $f$ est un monomorphisme propre.
\item[\emph{c)}] $f$ est propre et pour tout $y\in Y$, $f^{-1}(y)$ est
radiciel et géométriquement réduit sur $k(y)$ (c'est-à-dire vide ou
$k(y)$-isomorphe à $\operatorname{Spec}(k(y))$).
\end{enumerate}
\end{proposition}

Il est clair que \emph{a)} entraîne \emph{b)}. Pour voir que \emph{b)}
entraîne \emph{c)}, notons (I, 3.3.12) que pour tout $y\in Y$, le morphisme
$f^{-1}(y)\to\operatorname{Spec}(k(y))$ déduit de $f$ par extension de la base
est un monomorphisme, donc est injectif, et par suite $f^{-1}(y)$ est vide ou
réduit à un point, et en tout cas affine. De plus, si $A$ est l'anneau de
$f^{-1}(y)$, l'homomorphisme canonique $A\otimes_kA\to A$ est bijectif
(I, 5.3.8). Cela entraîne évidemment que $A=k$, sans quoi il y aurait un
élément $a\in A$ non dans $k$ et les images de $a\otimes1$ et $1\otimes a$
dans $A$ seraient toutes deux égales à $a$, alors que $a\otimes1\neq1\otimes a$
puisque $1$ et $a$ forment un système libre sur $k$.

Reste à prouver que \emph{c)} entraîne \emph{a)}. Il résulte tout d'abord de
(8.11.1) que $f$ est un morphisme \emph{fini}, donc
$X=\operatorname{Spec}(\mathcal A)$, où $\mathcal A$ est une
$\mathcal O_Y$-Algèbre \emph{finie}. Il suffit donc de prouver que
l'homomorphisme canonique $\mathcal O_Y\to\mathcal A$ est \emph{surjectif}
(II, 1.4.10), ou encore que pour tout $y\in Y$, l'homomorphisme
$\mathcal O_y\to\mathcal A_y$ est surjectif. Mais par hypothèse
